% GENERATED FILE. Edit canonical lessons, then run npm run generate:books.
% canonical-source-hash: fnv1a64:2584f0760f342bbe
% canonical-lessons: BN-C490-bichchhed, BN-C490-bidhoba, BN-C490-onath, BN-C490-dottok-neoya, BN-C490-obosor-neoya

\chapter{Separation and Loss}
\label{ch:bn-separation-and-loss}
\hlchaptermodality{490}

\begin{chapteropening}
\textbf{By the end of this chapter:} \emph{I can talk about separation, widowhood, adoption and retirement.}

Complete the last lesson of chapter 490: I can talk about separation, widowhood, adoption and retirement.
\end{chapteropening}

\section[\texorpdfstring{\bn{বিচ্ছেদ} (bichchhed)}{বিচ্ছেদ (bichchhed)}]{\bn{বিচ্ছেদ} (bichchhed) — a separation; a break-up}

\label{lesson:BN-C490-bichchhed}



\begin{quote}
Before the new one: say the Bengali for a blessing, then the Bengali for to touch an elder's feet.
\end{quote}

\subsection*{You'll want to know: \bn{বিচ্ছেদ}}
\textbf{\bn{বিচ্ছেদ}} — \emph{bichchhed} — \textquotedblleft{}a separation; a break-up\textquotedblright{}.

\begin{quote}
\bn{বিয়ের} \bn{দশ} \bn{বছর} \bn{পরে} \bn{মেসো} \bn{আর} \bn{মাসির} \bn{বিচ্ছেদ} \bn{হল।}
\end{quote}

\begin{quote}
\emph{biyer dôsh bôchhôr pôre meso ār māsir bichchhed hôlô.}
\end{quote}

\begin{quote}
Ten years after their wedding, my aunt and her husband separated.
\end{quote}

\begin{cousinweb}
\textbf{\bn{বিচ্ছেদ}} (\emph{bichchhed}) is a tatsama, Sanskrit \emph{viccheda}, \textquotedblleft{}cutting asunder, separation\textquotedblright{} (Monier-Williams).
\end{cousinweb}

\begin{grammarlens}[title={What you've built}]
One more word for describing what you lived and felt.
\end{grammarlens}

\subsection*{Your turn}
\begin{itemize}
\raggedright
  \item \emph{Say it:} \emph{bichchhed}
  \item \emph{Say it:} \emph{bichchhed}, once more
  \item \emph{Say it:} \emph{biyer dôsh bôchhôr pôre meso ār māsir bichchhed hôlô.} — Ten years after their wedding, my aunt and her husband separated.
  \item \emph{Recall:} say the Bengali for a blessing, then the Bengali for to touch an elder's feet, then say \emph{bichchhed} again
\end{itemize}

\begin{tcolorbox}[breakable,colback=teal!4,colframe=teal!35!black,title={Before you move on}]
What does \bn{বিচ্ছেদ} mean? (\textquotedblleft{}A separation; a break-up\textquotedblright{}.) Say it once more.
\end{tcolorbox}
\section[\texorpdfstring{\bn{বিধবা} (bidhôbā)}{বিধবা (bidhôbā)}]{\bn{বিধবা} (bidhôbā) — a widow}

\label{lesson:BN-C490-bidhoba}



\begin{quote}
Before the new one: say the Bengali for to touch an elder's feet, then the Bengali for a separation.
\end{quote}

\subsection*{You'll want to know: \bn{বিধবা}}
\textbf{\bn{বিধবা}} — \emph{bidhôbā} — \textquotedblleft{}a widow\textquotedblright{}.

\begin{quote}
\bn{যুদ্ধে} \bn{স্বামীর} \bn{মৃত্যুর} \bn{পরে} \bn{তিনি} \bn{বিধবা} \bn{হলেন।}
\end{quote}

\begin{quote}
\emph{juddhe svāmīr mrityur pôre tini bidhôbā hôlen.}
\end{quote}

\begin{quote}
After her husband died in the war, she was widowed.
\end{quote}

\begin{cousinweb}
\textbf{\bn{বিধবা}} (\emph{bidhôbā}) is a tatsama, Sanskrit \emph{vidhavā}, \textquotedblleft{}a husbandless woman, a widow\textquotedblright{} (Monier-Williams).
\end{cousinweb}

\begin{grammarlens}[title={What you've built}]
One more word for describing what you lived and felt.
\end{grammarlens}

\subsection*{Your turn}
\begin{itemize}
\raggedright
  \item \emph{Say it:} \emph{bidhôbā}
  \item \emph{Say it:} \emph{bidhôbā}, once more
  \item \emph{Say it:} \emph{juddhe svāmīr mrityur pôre tini bidhôbā hôlen.} — After her husband died in the war, she was widowed.
  \item \emph{Recall:} say the Bengali for to touch an elder's feet, then the Bengali for a separation, then say \emph{bidhôbā} again
\end{itemize}

\begin{tcolorbox}[breakable,colback=teal!4,colframe=teal!35!black,title={Before you move on}]
What does \bn{বিধবা} mean? (\textquotedblleft{}A widow\textquotedblright{}.) Say it once more.
\end{tcolorbox}
\section[\texorpdfstring{\bn{অনাথ} (ônāth)}{অনাথ (ônāth)}]{\bn{অনাথ} (ônāth) — an orphan; orphaned}

\label{lesson:BN-C490-onath}



\begin{quote}
Before the new one: say the Bengali for a separation, then the Bengali for a widow.
\end{quote}

\subsection*{You'll want to know: \bn{অনাথ}}
\textbf{\bn{অনাথ}} — \emph{ônāth} — \textquotedblleft{}an orphan; orphaned\textquotedblright{}.

\begin{quote}
\bn{বন্যায়} \bn{বাবাকে} \bn{হারিয়ে} \bn{ছেলেটা} \bn{অনাথ} \bn{হল।}
\end{quote}

\begin{quote}
\emph{bônyāy bābāke hāriye chheleṭā ônāth hôlô.}
\end{quote}

\begin{quote}
Losing his father in the flood, the boy was orphaned.
\end{quote}

\begin{cousinweb}
\textbf{\bn{অনাথ}} (\emph{ônāth}) is a tatsama, Sanskrit \emph{anātha}, \textquotedblleft{}having no master or protector\textquotedblright{} (Monier-Williams).
\end{cousinweb}

\begin{grammarlens}[title={What you've built}]
One more word for describing what you lived and felt.
\end{grammarlens}

\subsection*{Your turn}
\begin{itemize}
\raggedright
  \item \emph{Say it:} \emph{ônāth}
  \item \emph{Say it:} \emph{ônāth}, once more
  \item \emph{Say it:} \emph{bônyāy bābāke hāriye chheleṭā ônāth hôlô.} — Losing his father in the flood, the boy was orphaned.
  \item \emph{Recall:} say the Bengali for a separation, then the Bengali for a widow, then say \emph{ônāth} again
\end{itemize}

\begin{tcolorbox}[breakable,colback=teal!4,colframe=teal!35!black,title={Before you move on}]
What does \bn{অনাথ} mean? (\textquotedblleft{}An orphan; orphaned\textquotedblright{}.) Say it once more.
\end{tcolorbox}
\section[\texorpdfstring{\bn{দত্তক} \bn{নেওয়া} (dôttôk neoyā)}{দত্তক নেওয়া (dôttôk neoyā)}]{\bn{দত্তক} \bn{নেওয়া} (dôttôk neoyā) — to adopt (a child)}

\label{lesson:BN-C490-dottok-neoya}



\begin{quote}
Before the new one: say the Bengali for a widow, then the Bengali for an orphan.
\end{quote}

\subsection*{You'll want to know: \bn{দত্তক} \bn{নেওয়া}}
\textbf{\bn{দত্তক} \bn{নেওয়া}} — \emph{dôttôk neoyā} — \textquotedblleft{}to adopt (a child)\textquotedblright{}.

\begin{quote}
\bn{কাকা} \bn{আর} \bn{কাকিমা} \bn{একটা} \bn{অনাথ} \bn{মেয়েকে} \bn{দত্তক} \bn{নিলেন।}
\end{quote}

\begin{quote}
\emph{kākā ār kākimā ekṭā ônāth meyeke dôttôk nilen.}
\end{quote}

\begin{quote}
My uncle and aunt adopted an orphan girl.
\end{quote}

Past (he/she …): \bn{তিনি} \bn{একটা} \bn{শিশুকে} \bn{দত্তক} \bn{নিলেন}, \emph{tini ekṭā shishuke dôttôk nilen}.

\begin{cousinweb}
\textbf{\bn{দত্তক} \bn{নেওয়া}} (\emph{dôttôk neoyā}) is Sanskrit \emph{dattaka}, used with \emph{putra}, son, for a son given in adoption (Monier-Williams), with \emph{neoyā}, \textquotedblleft{}to take\textquotedblright{} (CDIAL 6966).
\end{cousinweb}

\begin{grammarlens}[title={What you've built}]
One more word for describing what you lived and felt.
\end{grammarlens}

\subsection*{Your turn}
\begin{itemize}
\raggedright
  \item \emph{Say it:} \emph{dôttôk neoyā}
  \item \emph{Say it:} \emph{dôttôk neoyā}, once more
  \item \emph{Say it:} \emph{kākā ār kākimā ekṭā ônāth meyeke dôttôk nilen.} — My uncle and aunt adopted an orphan girl.
  \item \emph{Recall:} say the Bengali for a widow, then the Bengali for an orphan, then say \emph{dôttôk neoyā} again
\end{itemize}

\begin{tcolorbox}[breakable,colback=teal!4,colframe=teal!35!black,title={Before you move on}]
What does \bn{দত্তক} \bn{নেওয়া} mean? (\textquotedblleft{}To adopt (a child)\textquotedblright{}.) Say it once more.
\end{tcolorbox}
\section[\texorpdfstring{\bn{অবসর} \bn{নেওয়া} (ôbôsôr neoyā)}{অবসর নেওয়া (ôbôsôr neoyā)}]{\bn{অবসর} \bn{নেওয়া} (ôbôsôr neoyā) — to retire}

\label{lesson:BN-C490-obosor-neoya}



\begin{quote}
Before the new one: say the Bengali for an orphan, then the Bengali for to adopt.
\end{quote}

\subsection*{You'll want to know: \bn{অবসর} \bn{নেওয়া}}
\textbf{\bn{অবসর} \bn{নেওয়া}} — \emph{ôbôsôr neoyā} — \textquotedblleft{}to retire\textquotedblright{}.

\begin{quote}
\bn{বাবা} \bn{গত} \bn{বছর} \bn{চাকরি} \bn{থেকে} \bn{অবসর} \bn{নিলেন।}
\end{quote}

\begin{quote}
\emph{bābā gôtô bôchhôr chākri theke ôbôsôr nilen.}
\end{quote}

\begin{quote}
Father retired from his job last year.
\end{quote}

Past (he/she …): \bn{তিনি} \bn{পঞ্চাশ} \bn{বছরে} \bn{অবসর} \bn{নিলেন}, \emph{tini pônchāsh bôchhôre ôbôsôr nilen}.

\bn{অবসর} alone is free time: \bn{অবসরে}, in one's spare time.

\begin{cousinweb}
\textbf{\bn{অবসর} \bn{নেওয়া}} (\emph{ôbôsôr neoyā}) is Sanskrit \emph{avasara}, \textquotedblleft{}occasion, opportunity, leisure\textquotedblright{} (Monier-Williams), with \emph{neoyā}, \textquotedblleft{}to take\textquotedblright{} (CDIAL 6966): to take your leisure for good.
\end{cousinweb}

\begin{grammarlens}[title={What you've built}]
One more word for describing what you lived and felt.
\end{grammarlens}

\subsection*{Your turn}
\begin{itemize}
\raggedright
  \item \emph{Say it:} \emph{ôbôsôr neoyā}
  \item \emph{Say it:} \emph{ôbôsôr neoyā}, once more
  \item \emph{Say it:} \emph{bābā gôtô bôchhôr chākri theke ôbôsôr nilen.} — Father retired from his job last year.
  \item \emph{Recall:} say the Bengali for an orphan, then the Bengali for to adopt, then say \emph{ôbôsôr neoyā} again
\end{itemize}

\begin{tcolorbox}[breakable,colback=teal!4,colframe=teal!35!black,title={Before you move on}]
What does \bn{অবসর} \bn{নেওয়া} mean? (\textquotedblleft{}To retire\textquotedblright{}.) Say it once more.
\end{tcolorbox}
